\section{结论}

本文针对自主水下航行器在粗分辨率海流图下执行强流绕障时的在线元决策问题，提出了基于海洋大数据标定的合理反例检验方法，从边缘智能与物理海洋学交叉的视角阐明了一条基本原理：机载计算的职责是消除求解器近似带来的数值模糊，从而忠实读出已有海流图所蕴含的倾向；而当多年海洋大数据所界定的合理流场集合内存在足以翻转航路优劣并造成严重航时损失的物理反例时，环境本身处于不可由算力消解的多解状态，航行器唯有通过低代价的大尺度观测引入新信息，或在观测昂贵时直接倒向最坏合理后悔更小的保守航路，才能规避灾难性的逆流风险。在台湾以东至东海黑潮区域六年再分析海流共1333个强流任务上的系统实验表明，将机载航线求解算力提高64倍仅能修复少部分粗算误差，而合理反例的最坏后悔则将高达67\%的地图总损失精准锁定在排名前20\%的高危任务中，使门控保守策略在不消耗额外观测时间的条件下将平均航时损失显著削减28\%，并清晰界定了原位测流的局限与上浮取新图的代价阈值。未来的工作将把合理反例检验从定常快照扩展至随时间连续演变的四维流场预报集合，并结合搭载声学多普勒海流剖面仪的自主水下航行器开展海上现场试验。
